%This is a experiment example of ZhengXiaoyang's experiment report template

\documentclass[UTF8]{ctexart}
 
\usepackage{amsmath}
\usepackage{cases}
\usepackage{cite}
\usepackage{xeCJK}
\usepackage{graphicx}
\usepackage[margin=1in]{geometry}
\geometry{a4paper}
\usepackage{fancyhdr}
\pagestyle{fancy}
\fancyhf{}

\graphicspath{{picture/}}


\title{实验全名}
\graphicspath{{picture/}}


\title{实验报告全名Title}
\author{郑晓旸}
\date{\today}
\pagenumbering{arabic}

\begin{document}
%这里是文件的开头
\fancyhead[L]{郑晓旸}
\fancyhead[C]{实验简称}
\fancyfoot[C]{\thepage}

\maketitle
\tableofcontents
\newpage

\section{实验目的}
    \

\section{实验仪器}


\section{实验原理}


\section{实验过程和数据分析}

\section{分析与讨论}

\subsection{误差分析}

\subsubsection{实验中的系统误差}


\subsubsection{实验中的偶然误差}
接线时可能有xxx情况，导致xxx。xx上的xx在某情况下有xx的问题存在，经反复调整后得以正常测量。

\subsection{实验后的思考}
可说明自己做本实验的总结、收获和体会，对实验中发现的问题提出自己的建议。

\newpage
%图一般很大，建议换页。
\section{原始数据}
\begin{center}
    Change the picture by yourself
    
    
    \includegraphics{picture/example.png}
\end{center}



\bibliographystyle{plain}
\bibliography{./template}  %bib文件名

\end{document}
